\DocumentMetadata{pdfversion=1.7,lang=zh-CN}
\documentclass[professional-degree,final-cover]{neudissertation}
\usepackage{algorithm}
\usepackage{algorithmic}
\AddToHook{env/table/begin}{\renewcommand{\arraystretch}{1.3}}
\newcommand{\sampletext}{这是一段用于观察垂直间距的测试正文。所有样例采用相同的字号与行距，以便比较正文与图形、题注、表格和算法之间的距离。测试内容仅用于排版，不属于论文正文。\par}
\newcommand{\panel}{\fbox{\rule{0pt}{24mm}\rule{90mm}{0pt}}}
\newcommand{\testfigure}[2]{\begin{figure}[#1]\centering\panel\neufigcaption{fig:#2}{用于观察间距的测试图形}{A test figure for vertical spacing}\end{figure}}
\newcommand{\testtable}[2]{\begin{table}[#1]\centering\neutabcaption{tab:#2}{用于观察间距的测试表格}{A test table for vertical spacing}\begin{tabular}{lcc}\toprule 项目 & 数值一 & 数值二\\\midrule 样例甲 & 12 & 24\\ 样例乙 & 18 & 36\\\bottomrule\end{tabular}\end{table}}
\newcommand{\testalgorithm}[1]{\begin{algorithm}[#1]\caption{用于观察间距的测试算法}\begin{algorithmic}[1]\REQUIRE 输入样本集合 $X$\ENSURE 输出结果集合 $Y$\STATE 初始化结果集合\FOR{每个输入样本}\STATE 计算并保存结果\ENDFOR\RETURN $Y$\end{algorithmic}\end{algorithm}}
\begin{document}
\mainmatter
\typeout{SPACING intextsep=\the\intextsep; textfloatsep=\the\textfloatsep; floatsep=\the\floatsep; baselineskip=\the\baselineskip; parskip=\the\parskip}
\chapter{垂直间距隔离测试}
\section{固定位置的图与表}
\sampletext
\testfigure{H}{fixed}
\sampletext
\testtable{H}{fixed}
\sampletext
\clearpage
\section{普通文中浮动与算法}
\sampletext
\testfigure{!htbp}{normal}
\sampletext
\testalgorithm{H}
\sampletext
\clearpage
\section{各级标题与正文}
\sampletext
\subsection{小节标题}
\sampletext
\subsubsection{更低一级的标题}
\sampletext
\paragraph{段落标题}
\sampletext
\section{连续标题的情形}
\subsection{标题之后直接出现下一级标题}
\sampletext
\clearpage
\section{图形之后直接出现标题}
\sampletext
\testfigure{H}{heading}
\subsection{紧接图题的小节标题}
\sampletext
\testalgorithm{H}
\subsection{紧接算法的小节标题}
\sampletext
\clearpage
\testfigure{t}{top}
\section{页顶浮动体与长题注}
\sampletext
\begin{figure}[H]\centering\panel
\neufigcaption{fig:long}{用于检查多行题注的较长中文图形标题，观察标题内部行距及其与后续正文的关系}{A deliberately long English figure caption used to inspect line wrapping and the distance between the caption block and the following body paragraph}
\end{figure}
\sampletext
\testtable{b}{bottom}
\sampletext
\clearpage
\section{节标题之后直接接固定图}
\testfigure{H}{section-before}
\sampletext
\subsection{小节标题之后直接接固定图}
\testfigure{H}{subsection-before}
\sampletext
\clearpage
\section{较低层级标题之后直接接图}
\sampletext
\subsubsection{三级标题之后直接接固定图}
\testfigure{H}{subsubsection-before}
\sampletext
\paragraph{段落标题之后直接接固定图}
\testfigure{H}{paragraph-before}
\sampletext
\clearpage
\section{节标题之后直接接普通浮动图}
\testfigure{!htbp}{section-before-normal}
\sampletext
\subsection{小节标题之后直接接普通浮动图}
\testfigure{!htbp}{subsection-before-normal}
\sampletext
\clearpage
\section{节标题之后直接接表格}
\testtable{H}{section-before}
\sampletext
\subsection{小节标题之后直接接算法}
\testalgorithm{H}
\sampletext
\subsection{表格之后直接接小节标题}
\testtable{H}{heading-after-table}
\subsection{紧接表格的小节标题}
\sampletext
\clearpage
\section{接近页末的标题与图形}
\sampletext\sampletext\sampletext\sampletext\sampletext\sampletext\sampletext
\sampletext\sampletext\sampletext\sampletext\sampletext
\subsection{标题与下方图形能否一起分页}
\testfigure{H}{page-break}
\sampletext
\end{document}
